\section{An alternative globalization by positive degree}
\label{sec:degree-alternative}

The covering argument in the main text uses local injectivity of the
graph projection.  Here we give a different sufficient condition:
openness, together with local minimality of its action branches, also
forces global injectivity.  The reason is that local minimality makes
the projection Jacobian nonnegative in coordinates on the subgradient
graph, while the homotopy from a short-time projection has degree one.
The area formula then rules out two overlapping open images.
This argument uses the same active logs and divided-action estimates
as the main proof; it does not require the global supporting theorem.

\subsection{Local minima and openness}

\begin{lemma}[Local quadratic growth]
\label{geo:local-growth}
Let $u$ be a potential and $0<t<1$.  Suppose that, for every
$(x,p)\in\Gamma_u$, one has $rp\in I(x)$ for $0<r<t$ and $tp$
is nonconjugate.  Put $z=\exp_x(tp)$ and let $\hat c$ be the smooth
action branch through $(x,tp)$.  Then there are $b,\rho>0$ such that
\[
 u(x')+\frac{\hat c(x',z)}t
 \geq u(x)+\frac{\hat c(x,z)}t+b\,d(x,x')^2
 \qquad\text{if }d(x,x')<\rho.
\]
In particular, $F_t:\Gamma_u\to M$ is open.
\end{lemma}

\begin{proof}
Choose a positive finite active representation
$p=\sum_i m_ip_i$, and put $E=\Span\{p_i-p\}$.
Fix $s$ with $t<s<1$.  Each $p_i$ is an original minimizing log,
so its shortened support at time $s$ is smooth near $x$.
In normal coordinates $h$ at $x$, these supports give
\[
 u(\exp_xh)-u(x)
 \geq\max_i\left\{\langle p_i,h\rangle
                 -\frac12H_s(p_i)(h,h)\right\}+o(|h|^2).
\]
Adding the selected branch of $\hat c/t$ yields
\begin{equation}
 \label{geo:growth-expansion}
 \begin{split}
 &u(\exp_xh)+\frac{\hat c(\exp_xh,z)}t
       -u(x)-\frac{\hat c(x,z)}t\\
 &\qquad\geq
 \max_i\left\{\langle p_i-p,h\rangle
                     +\frac12B_i(h,h)\right\}+o(|h|^2),
 \qquad B_i=H_t(p)-H_s(p_i).
 \end{split}
\end{equation}
The contracted active hulls lie in $I(x)$ by hypothesis, and their
limits at time $t$ are minimizing.  The branch at $tp$ is
nonconjugate.  Thus the limiting form of \cref{geo:slack} gives
\[
 \left.\sum_i m_iB_i\right|_{E^\perp}
 \geq\kappa_g(s-t)\Id.
\]
Only the individual active rays are evaluated at the later time $s$.

The vectors $p_i-p$ span $E$, have positive weights, and have weighted
mean zero.  Hence, if $E\ne\{0\}$, compactness of its unit sphere gives
\[
 \max_i\langle p_i-p,h\rangle\geq b_E|h_E|,
 \qquad h_E=\operatorname{proj}_Eh,
 \qquad b_E>0.
\]
For $E\ne\{0\}$, choose $B$ so that
$b_E B>1+\max_i\|B_i\|$.
When $|h_E|\geq B|h|^2$, this linear term dominates all negative
quadratic terms in \eqref{geo:growth-expansion} and leaves a positive
multiple of $|h|^2$.  When $|h_E|<B|h|^2$, take the weighted average
inside the maximum.  The linear term vanishes, the quadratic term on
$E^\perp$ is positive, and its mixed terms are $O(|h|^3)$.
The same averaged argument applies directly if $E=\{0\}$.
After shrinking the neighborhood, the claimed quadratic growth follows.

For openness, choose a small closed ball about $x$ inside this
neighborhood and inside the source domain of the branch.  For an
endpoint $z'$ near $z$, minimize
$u(\cdot)+\hat c(\cdot,z')/t$ on that ball.  The positive boundary
gap keeps each minimizer $x'$ in its interior.  As $z'\to z$, every
such minimizer tends to $x$, since quadratic growth makes $x$ the
unique minimizer on the ball at $z$.  Stationarity gives
\[
 p'=-\frac1tD_x\hat c(x',z')\in\partial u(x').
\]
First variation on the branch gives $\exp_{x'}(tp')=z'$ and
$(x',p')\to(x,p)$.  Thus every neighborhood of $(x,p)$ in the graph
projects onto a neighborhood of $z$, as required.
\end{proof}

\subsection{The sign of a graph projection}

The next lemma is conditional: its local-minimum assumption concerns
the terminal time $t$ only.  The required short-time properties follow
separately from the semiconvexity bounds and the short-time
homeomorphism already proved in the main text.

\begin{lemma}[Positive projection degree]
\label{geo:degree}
Let $u$ be a potential and $0<t<1$.  Assume that every $rp$, with
$(x,p)\in\Gamma_u$ and $0<r\leq t$, is nonconjugate.  At each
$(x,p)$, put $z=\exp_x(tp)$ and let $\hat c$ be the smooth action
branch determined by $tp$.  Suppose that $x$ is a local minimum of
\[
 x'\longmapsto u(x')+\frac{\hat c(x',z)}t,
\]
and suppose that $F_t:\Gamma_u\to M$ is open.
Then $F_t$ is a homeomorphism.  Its pole at every endpoint is the
unique global minimizing pole of $Q_tu$, and every $tp$ lies in $I(x)$.
\end{lemma}

\begin{proof}
Choose a sufficiently small $\eps\in(0,t)$ as in
\cref{geo:short-time}.  That result makes $F_\eps$ a homeomorphism
with locally Lipschitz inverse.  Since $F_\eps$ is the restriction
of a smooth map on the compact graph, these are bi-Lipschitz graph
charts.

We use Minty-type coordinates to compute the Jacobian sign.  In a
coordinate ball write $p$ as a covector and take $a>0$ small enough that
$h(x)=|x|^2/2+a u(x)$ is strongly convex.  The coordinates
\[
 q=x+ap
\]
parametrize the graph locally.  To see this, restrict $h$ to a small
closed ball around the reference pole and maximize $q\cdot x-h(x)$.
Strong convexity makes the maximizer unique and Lipschitz in $q$;
for $q$ near the reference value it is interior.  Thus $x(q)$ is
the gradient of the local convex conjugate, and
$p(q)=(q-x(q))/a$.  This parametrization is bi-Lipschitz.
For the classical convex antecedents, see the proximal parametrization
in \cite[Section~12]{Moreau1965} and the graph correspondence in
\cite[Proposition~1.1]{AA1999}; compare also
\cite[Section~2]{MPW2012} for related diagonal coordinates in optimal
transport.

Rademacher's theorem \cite[Chapter~2, Theorem~1.4]{Simon2018} gives,
at almost every point of a graph chart,
\begin{equation}
 \label{geo:minty-matrices}
 X=D_qx\geq0,
 \qquad P=D_qp=\frac{\Id-X}{a}.
\end{equation}
Both matrices are symmetric and commute.  Write
$\mathscr A=D^2_{xx}\hat c(x,F_t(x,p))$ for the coordinate source
Hessian of the selected unit-duration action and put $Z=D_qF_t$.
Differentiating $-D_x\hat c(x,F_t)=tp$ gives
\begin{equation}
 \label{geo:projection-matrix}
 Z=(-D^2_{xz}\hat c)^{-1}(tP+\mathscr A X).
\end{equation}

Local minimality at the terminal time implies
\begin{equation}
 \label{geo:sign-form}
 X(tP+\mathscr A X)\geq0.
\end{equation}
Here is a justification requiring only first derivatives of the graph.
Along a straight line in the chart, put
$x_r=x(q+r\xi)$ and $p_r=p(q+r\xi)$.  The semiconvex
subgradient inequalities at two parameter values imply, almost
everywhere on the line,
\[
 \frac{d}{dr}u(x_r)=\langle p_r,x'_r\rangle.
\]
At a differentiability point of $(x,p)$, the expansions
$x_r=x+rX\xi+o(r)$ and $p_r=p+rP\xi+o(r)$ give
\[
 \begin{split}
 u(x_r)-u(x)-\langle p,x_r-x\rangle
 &=\int_0^r\langle p_s-p,x'_s\rangle\,ds\\
 &=\frac{r^2}{2}\langle P\xi,X\xi\rangle+o(r^2).
 \end{split}
\]
For the last equality, integrate the term $sP\xi$ by parts.
The remaining integral is $o(r^2)$ because $x_s$ is Lipschitz and
the remainder in $p_s$ is $o(s)$.  No second derivative of the pole
curve is used.
Rewrite this as an expansion of $u(x_r)-u(x)$, multiply by $t$,
and add the expansion of
$\hat c(x_r,z)-\hat c(x,z)$ with $z=F_t(x,p)$ fixed.
Its linear term is $-t\langle p,x_r-x\rangle$, and its
quadratic term is $r^2\langle\mathscr A X\xi,X\xi\rangle/2$.
The cancellation and the local-minimum inequality prove
\eqref{geo:sign-form}.

Split the matrices into $\ker X$ and its orthogonal complement.
Then $tP+\mathscr A X$ is block triangular, with kernel block
$(t/a)\Id$.  On the complementary block, multiplication by the
positive definite matrix $X$ makes it symmetric positive semidefinite
by \eqref{geo:sign-form}.  Consequently
\[
 \det(tP+\mathscr A X)\geq0,
\]
with positive sign whenever this matrix is nonsingular.
The other factor in \eqref{geo:projection-matrix} has positive
determinant in compatible oriented charts.  Indeed,
$-D^2_{xz}\hat c$ is the inverse exponential differential followed
by the metric identification.  Its sign starts positive at time zero
and cannot change along the ray before a conjugate time.  The
nonconjugacy assumption supplies this sign through time $t$.

Assume first that $M$ is oriented and transfer the projections to $M$:
\[
 \mathcal H_r=F_r\circ F_\eps^{-1}:M\to M,
 \qquad \eps\leq r\leq t.
\]
These maps are Lipschitz and form a homotopy from the identity to
$\mathcal H_t$.  At the small time $\eps$, semiconvexity and the
positive cost Hessian make each stationary graph branch a local
minimum; equivalently, $u+c_\eps(\cdot,z)$ is strongly convex on
the relevant coordinate ball.  The preceding sign calculation thus
applies to $F_\eps$ independently of the terminal-time assumption.
Its bi-Lipschitz differential is nonsingular almost everywhere, so
the computed sign is positive.  The chain rule transfers the
nonnegative sign of $DF_t$ to the signed Riemannian Jacobian
$J_{\mathcal H_t}$.  Bi-Lipschitz changes of coordinates preserve
the exceptional null sets, and hence
$J_{\mathcal H_t}\geq0$ almost everywhere.

For completeness, the signed degree identity can be justified on
the smooth manifold without imposing a smooth structure on the
graph.  Smoothly approximate $\mathcal H_t$ uniformly and strongly
in $W^{1,n}$, using local mollification and a smooth tubular
projection after embedding the compact target.  Uniformly close
approximations are homotopic to $\mathcal H_t$, hence to the
identity.  Stokes' theorem gives
$\int_M h^*\omega=\int_M\omega$ for every smooth top form $\omega$
and each such smooth approximation $h$.  Strong derivative
convergence passes this identity to $\mathcal H_t$: the pullback
is multilinear of degree $n$ in the derivative, and its coefficients
converge uniformly.  This identity also proves surjectivity.  A
missed point would give, by compactness, a missed open ball, and a
top form supported in that ball with nonzero integral would
contradict the identity.

Taking $\omega=d\vol$ and applying the Lipschitz area formula
\cite[Chapter~2, Theorem~3.3, Remark~3.4, and (4.19)--(4.20)]{Simon2018} in manifold
charts gives
\[
 \int_M J_{\mathcal H_t}\,d\vol=\vol(M),
 \qquad
 \int_M N(y,\mathcal H_t)\,d\vol(y)
   =\int_M|J_{\mathcal H_t}|\,d\vol=\vol(M),
\]
where $N$ counts preimages and may be infinite.  Surjectivity gives
$N\geq1$ everywhere, so $N=1$ almost everywhere.  If there were
two distinct preimages, disjoint neighborhoods of them would have
open images meeting in a nonempty open set, since $F_t$ and hence
$\mathcal H_t$ are open.  The multiplicity would be at least two
on that set, contradicting the integral equality.  Thus
$\mathcal H_t$ is injective.  No discreteness of fibers is assumed
in this argument.

If $M$ is nonorientable, lift the homotopy to its orientation double
cover, starting from the identity.  Uniqueness of homotopy lifting
shows that every lift commutes with the deck transformation.
The lifted maps are Lipschitz, and the terminal lift is open because
the covering maps are local diffeomorphisms.  Its local Jacobian
signs are the same as above: the lifted homotopy selects the target
orientation transported along the geodesic used in the mixed-action
calculation.  The oriented argument therefore makes the terminal
lift injective.  Two distinct preimages downstairs could be lifted,
applying a deck transformation to one of them if needed, to have
the same lifted image.  Deck equivariance would contradict
injectivity upstairs.  Thus $F_t$ is injective in either case;
surjectivity and compactness make it a homeomorphism.

Finally, take any global minimizing pole of $Q_tu$ at an endpoint.
\Cref{geo:minimizer-capture} places that pole and every
one of its minimizing logs divided by $t$ in the graph fiber.
Its uniqueness identifies the pole with the specified graph pole
and identifies every minimizing log with $tp$.  Thus $tp$ is the
unique minimizing velocity to its endpoint.  It is nonconjugate
by hypothesis, so \cref{lem:interior-characterization} places it in
$I(x)$ and proves the remaining assertions.
\end{proof}
